% 【本文件写什么】impl 实现层：platform-console
% - 定位：经 platform 固件/控制台的输出后端。
% - crate 路径：os/components/wateros-runtime/runtime-console/console-impl/impl-platform-console/src/
% - 如何实现对应 api-v0 trait：关键类型、状态机、数据结构
% - 算法或硬件交互流程，可用伪代码/流程图
% - 本 impl 启用的 Cargo [features] 与 arch feature，impl-riscv64 等
% - 与其它 impl 变体的差异、切换 feature 时的影响
% - 性能/内存假设与已知限制
% 【事实来源】
% - 上述 crate 源码与 Cargo.toml
% - os/feature-tree.txt 中指向本 impl 的 feature 链

\subsubsection{impl-platform-console}

\code{PlatformConsoleHandle}实现\code{Console}与\code{fmt::Write}：\code{write\_str}把UTF-8切片按字节经\code{platform::console::console\_write\_a\_buffer}下发，错误时\code{unwrap}，引导阶段视为致命失败。\code{platform\_console\_write\_a\_byte}同样委托平台门面。具体硬件路径，OpenSBI或MMIO UART由当前\code{wateros-platform} board feature决定，本crate不选择后端。

\begin{rustcode}
impl Write for PlatformConsoleHandle {
    fn write_str(&mut self, s: &str) -> fmt::Result {
        platform_console_write_a_buffer(s.as_bytes());
        Ok(())
    }
}
\end{rustcode}
